% Source: https://artofproblemsolving.com/community/c3h2486182_fraction_to_repeating_decimal
% Stable identity: aops-topic-2486182-q01
\textbf{(a)} Since $d\ge1$ we have $1\le N\le10^k-1$, so $N$ has at most $k$ digits; write it with exactly $k$ digits as $a_1a_2\dots a_k$. Then
\[\dfrac1d=\dfrac{N}{10^k-1}=\dfrac{N}{10^k}\cdot\dfrac{1}{1-10^{-k}}=N\left(10^{-k}+10^{-2k}+10^{-3k}+\cdots\right).\]
The term $N\cdot10^{-k}$ is $0.a_1a_2\dots a_k$, the term $N\cdot10^{-2k}$ places the same block in decimal places $k+1$ to $2k$, and so on. The blocks occupy disjoint sets of decimal places, so no carrying occurs and
\[\dfrac1d=0.\overline{a_1a_2\dots a_k}.\]

\textbf{(b)} Here $k=3$ and $N=\dfrac{999}{37}=27$, written as $027$. So
\[\dfrac{1}{37}=\boxed{0.\overline{027}}.\]
Also $\dfrac{5}{37}=\dfrac{135}{999}$, and the same argument with numerator $135$ gives $\dfrac{5}{37}=\boxed{0.\overline{135}}$.

\textbf{(c)} Here $k=6$ and $N=47619$, written with six digits as $047619$. So
\[\dfrac{1}{21}=\boxed{0.\overline{047619}}.\]
